\chapter*{Resumen}

En este Trabajo he abordado el estudio, evaluación y optimización de la navegación autónoma en robots móviles mediante técnicas de SLAM (Simultaneous Localization and Mapping). El objetivo principal ha sido analizar el comportamiento de distintos métodos de localización, mapeado y planificación de trayectorias en entornos simulados, con el fin de mejorar la precisión, robustez y eficiencia del sistema de navegación.

Para ello, he trabajado con el robot TurtleBot4 integrado en el ecosistema ROS 2 Humble, sobre el cual he implementado y probado diferentes configuraciones del stack de navegación.

El trabajo se ha desarrollado primero en entornos simulados controlados y posteriormente en entornos reales con el robot físico, lo que me ha permitido validar de forma progresiva el comportamiento de los algoritmos y transitar desde condiciones idealizadas hasta escenarios más realistas.

Esta transición ha sido especialmente relevante, ya que en el entorno real aparecen factores no presentes en simulación, como ruido en los sensores, pequeñas imprecisiones en el movimiento del robot o variaciones no modeladas del entorno. Gracias a esta combinación de simulación y pruebas reales, he podido reproducir escenarios con distintos niveles de complejidad y condiciones de operación, lo que ha facilitado una evaluación más completa del sistema de navegación y de su capacidad de generalización.

Este trabajo me ha permitido profundizar en el funcionamiento de los sistemas SLAM y comprender el impacto de sus configuraciones, contribuyendo al desarrollo de soluciones más eficientes en el ámbito de la robótica móvil.
